\documentclass{article}
\usepackage{graphicx} % Required for inserting images
\usepackage{xfrac}
\usepackage{amsmath}
\usepackage{amsthm}
\usepackage{amssymb}
\newtheorem{theorem}{Theorem}

\title{Strong Hasse Principle}
\author{Haim Lavi, 038712105}
\date{August 2026}

\begin{document}

\maketitle

\begin{theorem}
Let $(V,q)$ be an $n$-dimensional quadratic space over a field $K$, then a vector $v\in{K}^n$ is isotropic (that is, $v$ is a nontrivial solution of $q(v)=0$) for $K=\mathbb{Q}$ if and only if $v$ is isotropic for $K=\mathbb{Q}_p$ for all primes $p$ (including $\mathbb{R}=\mathbb{Q}_\infty$).
\end{theorem}
\begin{proof}
The only if direction is trivial, since $\mathbb{Q}\subset\mathbb{R}$ and $\mathbb{Q}\subset\mathbb{Q}_p$ for every prime $p$ (the image of every $r\in\mathbb{Q}$ in either of its completion fields is the equivalence class of the constant Cauchy sequence $r,r,\dots,r,\dots$), but the if direction is complicated and requires separate proofs for several cases of $n$.
\section*{$n=2$}
Consider the quadratic form $q=a_1x_1^2+a_2x_2^2$, and assume that $a_1=1$ (this will be the assumption for this case and for the next one)\footnote{otherwise divide $q$ by $a_1$, then for all $v\in{K}^2$ we have $q(v)=0\iff\frac{1}{a_1}q(v)=0$, and we take $q$ to be $\frac{1}{a_1}q$.}.\\
If $v=(v_1,v_2)\in\mathbb{R}^2$ is a solution of $q$, then $q(v)=v_1^2+a_2v_2^2=0\iff{a_2=-(\frac{v_1}{v_2}})^2$, thus $-a_2>0$ (because it is a square in $\mathbb{R}$)\footnote{This is true for $\mathbb{R}$ because it has an ordering while $\mathbb{Q}_p$, for any prime $p$, does not (as follows from the Artin–Schreier theorem).
We can show this explicitly by considering the polynomial $f(x)=x^2-(1-p)$ (hence $f'(x)=2x$) for some prime $p$, then $f(1)=1-1+p\equiv0(\mod{p})$, and $f'(1)=2\neq0$, thus by Hensel’s lemma, there exists $a$ such that $f(a)=0$, which means that $a^2=1-p$, but then $b=a^2+1^2(p-2)$ hence $b+1=0\implies{b}=-1$.
}, so we can reformat $q$ as $q=x_1^2-a_2x_2^2$, where $a_2>0$.\\
\newline
There exists also a solution $v\in\mathbb{Q}_p^2$, for every prime $p$, hence $a_2=(\frac{v_1}{v_2})^2$ is a square in $\mathbb{Q}_p$, but, denote $r=\frac{v_1}{v_2}$, $a_2$ can be factored as \[
a_2=\prod_pp^{v_p(a_2)}=\prod_pp^{v_p(r^2)}=\prod_pp^{2v_p(r)}=\prod_p\left(p^{v_p(r)}\right)^2,
\]
which means that $a_2$ is a product of squared integers, hence $a_2=a_2'^2$, for some integer $a_2'\in\mathbb{Z}$, thus, taking $v_2=\frac{v_1}{a_2'}$ shows that $v'=(v_1,\frac{v_1}{a_2'})\in\mathbb{Q}^2$ yields $q(v')=v_1^2-a_2(\frac{v_1}{a_2'})^2=v_1^2-a_2'^2(\frac{v_1}{a_2'})^2=v_1^2-v_1^2=0$.
\section*{$n=3$}
Here $q=x_1^2+a_2x_2^2+a_3x_3^2$.\\
We establish first that $a_2$ and $a_3$ are square free integers\footnote{otherwise, suppose that $a_2$ is not square free, that is $a_2=bc$ where $b$ is square free and $c=\pm{p}_1^2p_2^2\cdots{p}_k^2=\pm{c}'^2$, where $c'={p}_1p_2\cdots{p}_k$ for $k$ distinct primes, but then $q=x_1^2-(bc)x_2^2-a_3x_3^2=x_1^2-b(c'x_2)^2-a_3x_3^2$, hence we transition $c'x_2\rightarrow{x}_2'$ to obtain the quadratic form $q'=x_1^2-bx_2'^2-a_3x_3^2$, and then for all $v=(v_1,v_2,v_3)\in{K}^3$ such that $q(v)=0$, we have $q'(v)=v_1^2-b(c'v_2)^2-a_3v_3^2=v_1^2-a_2v_2^2-a_3v_3^2=0$, and we can take $q'$ to be $q$. Same holds for $a_3$.}, but this means that for each prime $p$, the $p$-adic valuations $v_p(a_2),v_p(a_3)\in\{0,1\}$ ($1$ if $p$ is in the factorization, $0$ if it is not).\\
Without loss of generality we assume that $|a_2|\leq|a_3|$. Since $a_2$ and $a_3$ are integers, then their absolute values in $\mathbb{Q}$ are positive integers\footnote{while this assertion is true for (square free integers of) any field $K$, it is relevant here only for $\mathbb{Q}$, because this is the  target of this proof.}, hence their sum is a positive integer, and we denote $m=|a_2|+|a_3|\geq2$.\\
We start with $m=2$, that is $a_2=\pm1$ and $a_3=\pm1$, which means that \\$q=x_1^2\pm{x}_2^2\pm{x}_3^2$, but $q(v)=0$ over $\mathbb{R}$, so $a_2$ and $a_3$ cannot be both $1$, and it is immediate to find nonzero vectors $v=(v_1,v_2,v_3)\in\mathbb{Q}^3$ such that $q(v)=0$ (for instance, $(1,1,0),(1,0,1)\in\mathbb{Q}^3$ are both nontrivial solutions of $q=x_1^2-x_2^2-x_3^2$), thus we establish the induction base.\\
\newline
Now assume that $m>2$, which means that $|a_3|\geq{2}$ (because we assumed that $|a_3|\geq|a_2|$). We first notice that $q$ is of the form $q=x_1^2-a_2x_2^2-a_3x_3^2$ (or can be rearranged into this form)\footnote{because, as in the above case, $a_2$ and $a_3$ cannot be both positive and the coefficient of $x_1$ is always $1$, hence we can always reshuffle $x_1^2$, $a_2x_2^2$ and $a_3x_3^2$ to obtain the desired form.}.\\ 
Since $a_3$ is a square free integer, that is $a_3=p_1p_2\cdots{p}_k$ for $k$ distinct primes, then for an arbitrary prime $p_i$ of these primes, denote $p'=p_1p_2\cdots{p}_{i-1}p_{i+1}\cdots{p}_k$, thus $a_3=\pm{p}_ip'$, and there exists $v=(v_1,v_2,v_3)\in\mathbb{Q}_{p_i}^3$ such that $q(v)=v_1^2-a_2v_2^2-a_3v_3^2=0$, because $q$ has nontrivial solutions for all primes $p$, which means that ${v}_1^2-a_2v_2^2=a_3v_3^2=p_ip'v_3^2\equiv0(\mod{p_i})$. We assume that $v$ is primitive, that is, the common divisor of $v_1,v_2,v_3$ is $d=1$, otherwise $d>1$, and we can divide by it to obtain $u=\frac{v}{d}$, but $u_1^2-a_2u_2^2-a_3u_3^2=\frac{1}{d^2}(v_1^2-a_2v_2^2-a_3v_3^2)=0$, hence $q(v)=0\iff{q}(u)=0$, and we can take $u$ to be $v$.\\
\newline
Now we want to establish that $a_2$ is a square modulo $p_i$.
Suppose that $a_2\equiv0(\mod{p_i})$, that is $a_2=p_ia_2'$ for some $a_2'\in\mathbb{Z}_{p_i}$, then $a_2$ is trivially a square modulo $p_i$, because then $a_2$ and $0$ are in the same congruence class modulo $p_i$ (since $0=p_i0$, and $0$ is a square of itself), hence we assume that $a_2\not\equiv0(\mod{p_i})$ and move to show that it is a (nonzero) square modulo $p_i$.\\
\newline
We now look at $v_2$. Suppose that $v_2\equiv0(\mod{p_i})$, that is $v_2=p_iv_2'$ for some $v_2'\in\mathbb{Z}_p$, then recalling that
$a_3v_3^2=p_ip'v_3^2\equiv0(\mod{p_i})$, we have $v_1^2=a_2v_2^2+a_3v_3^2=a_2(p_iv_2')^2+p_ip'v_3^2=p_i(a_2p_iv_2'+p'v_3^2)\implies{p_i}|v_1^2\implies{p_i}|v_1$, because $p_i$ is prime, hence $v_1=p_iv_1'$ for some $v_1'\in\mathbb{Z}_p$, but then $a_3v_3^2=v_1^2-a_2v_2^2=(p_iv_1')^2-a_2(p_iv_2')^2\implies{p_i}^2|a_3v_3^2$, which means that $v_{p_i}(a_3v_3^2)=v_{p_i}(a_3)+v_{p_i}(v_3^2)=v_{p_i}(a_3)+2v_{p_i}(v_3)\geq2$, but $v_{p_i}(a_3)=v_p(p_ip')=1$, but $v_{p_i}(v_3)\geq1\iff{p_i}|v_3$, thus $p_i$ is a common divisor of $v_1,v_2,v_3$, in contradiction to $(v_1,v_2,v_3)$ being primitive, from which we conclude that $v_2\not\equiv0(\mod{p}_i)\implies{v}_2\in\mathbb{Z}_p^{\ast}\implies{v}_2^2\in\mathbb{Z}_p^{\ast}$ (that is, $v_2$ and $v_2^2$ are both $p$-adic invertible integers)\footnote{Let $a\in\mathbb{Z}_p$, then $a\in\mathbb{Z}_p^{\ast}\iff{v}_p(a)=0$, because $v_p(a)=0$ means that $a=a_0+pa'$ for some $1\leq{a}_0\leq(p-1)$, hence we consider the polynomial $f(x)=ax-1$, and by Euler’s Theorem there exists an integer $x_0$ such that $x_0a_0\equiv1(\mod{p})$ that is $x_0a_0=1+pa_0'$ for some integer $a_0'$, and then $f(x_0)=ax_0-1=a_0x_0+a'px_0-1=1+pa_0'+a'px_0-1=p(a_0'+a'x)\equiv0(\mod{p})$, and since $f'(x_0)=a\neq0$, then by Hensel’s lemma there exists $b$, such that $f(b)=ab-1=0$, which mean that $b=\frac{1}{a}$ (the opposite direction is almost immediate, and so is claim regarding the square of an invertible $p$-adic integer).}.
We rearrange the equation $v_1^2-a_2v_2^2=p_ip'v_3^2$ as $a_2v_2^2=v_1^2-p_ip'v_3^2$, but we saw that $v_2^2\in\mathbb{Z}_p^{\ast}\implies{a}_2=(v_2^2)^{-1}(v_1^2-p_ip'v_3^2)=(v_2^{-1})^2v_1^2-(v_2^{-1})^2p_ip'v_3^2\implies(v_2^{-1}v_1)^2=a_2+(v_2^{-1})^2p_ip'v_3^2\implies(v_2^{-1}v_1)^2\equiv{a}_2(\mod{p_i})$, thus $a_2$ is again a square modulo $p_i$.\\
\newline
By the Chinese Remainder Theorem we know that\footnote{The canonical map \[f:\sfrac{\mathbb{Z}}{a_3\mathbb{Z}}\rightarrow\prod_{j=1}^k\sfrac{\mathbb{Z}}{p_j\mathbb{Z}}\]
defined by $f(x+a_3\mathbb{Z})=(x+p_1\mathbb{Z},x+p_2\mathbb{Z},\dots,x+p_k\mathbb{Z})$ can be verified as preserving the ring operations, that is \[f\left((x+a_3\mathbb{Z})+(y+a_3\mathbb{Z})\right)=f\left((x+y)+a_3\mathbb{Z}\right)=\left((x+y)+p_1\mathbb{Z},\dots,(x+y)+p_k\mathbb{Z}\right)=\]\[=(x+p_1\mathbb{Z},\dots,x+p_k\mathbb{Z})+(y+p_1\mathbb{Z},\dots,y+p_k\mathbb{Z})=f(x+a_3\mathbb{Z})+f(y+a_3\mathbb{Z})\]
\[f\left((x+a_3\mathbb{Z})(y+a_3\mathbb{Z})\right)=f\left(xy+a_3\mathbb{Z}\right)=(xy+p_1\mathbb{Z},\dots,xy+p_k\mathbb{Z})=\]\[=(x+p_1\mathbb{Z},\dots,x+p_k\mathbb{Z})(y+p_1\mathbb{Z},\dots,y+p_k\mathbb{Z})=f(x+a_3\mathbb{Z})f(y+a_3\mathbb{Z})\]
and $f$ is a bijection as follows from the (the existence of a unique solution modulo $a_3$ to a system of congruences modulo $p_1,p_2,\dots,p_k$ as given by the) Chinese Remainder Theorem.
} \[\sfrac{\mathbb{Z}}{a_3\mathbb{Z}}=\sfrac{\mathbb{Z}}{(p_1p_2\cdots{p}_k)\mathbb{Z}}\cong\prod_{j=1}^k\sfrac{\mathbb{Z}}{p_j\mathbb{Z}},\]
and since all the above applies to every prime $p_j$ in the $k$ prime factors of $a_3$, because $q$ has a solution for all primes, then $a_2$ is a square modulo $p_j$ for every $1\leq{j}\leq{k}$, thus $a_2$ is also a square modulo $a_3$, that is $z^2=a_2+a_3y$ for some $y,z\in\mathbb{Z}$, and we can choose\footnote{otherwise, suppose that $|z|>\frac{|a_3|}{2}\iff|z-a_3|<\frac{|a_3|}{2}$, then $(a_3-z)^2=a_3-2a_3z+z^2=a_3(1-2z)+a_2+a_3y=a_2+a_3y'$ for some $y'$, thus $a_3-z$ is also a square modulo $a_3$, and we take $z-a_3$ to be $z$.} an integer $z$ such that $|z|\leq\frac{|a_3|}{2}$, and we observe that since $z^2-a_2=a_3y$, the equation $q'=x_1^2-a_2x_2^2-a_3yx_3^2$ has a nontrivial solution $v=(z,1,1)\in\mathbb{Q}_p^3$, because $q'(v)=z^2-a_2-a_3y=0$, thus $v$ is a solution\footnote{The fact that $q'$ has a nontrivial solution when we have the identity $z^2-a_2=a_3y$ can be generalized into the following proposition:\\
Let $q=x_1^2-a_2x_2^2-a_3x_3^2$, where $a_2,a_3\in{K}$, then $q$ has a nontrivial solution if and only if we have the identity $z^2-a_2=a_3$, which means that $a_3$ is a norm in the extension field $K[\sqrt{a_2}]$.}.\\
\newline
Suppose that $y=0\iff{z}^2=a_2$, that is $a_2$ is a square in $\mathbb{Z}$, then any vector of the form $v=(v_1,\frac{v_1}{z},0)\in\mathbb{Q}^3$ is a solution of $q$, because $q(v)=v_1^2-z^2(\frac{v_1}{z})^2-a_30^2=0$.\\
Otherwise, if $y\neq0$, then $|y|=|\frac{z^2-a_2}{a_3}|\leq|\frac{z^2}{a_3}|+|\frac{a_2}{a_3}|$, but we assumed that\\ $|z|\leq\frac{|a_3|}{2}$ and that $|a_2|\leq|a_3|$, hence $|y|\leq|\frac{z^2}{a_3}|+|\frac{a_2}{a_3}|\leq|\frac{1}{a_3}(\frac{a_3}{2})^2|+1=\frac{|a_3|}{4}+1$, and for every $|a_3|\geq2$, we have that $|y|\leq\frac{|a_3|}{4}+1<|a_3|$. \\
Now we claim the following:\\
There exists a bijection between the set of rational solutions of the (original) equation $q=x_1^2-a_2x^2-a_3x_3^2=0$ and that of  $q'=x_1^2-a_2x_2^2-yx_3^2=0$ (a full proof is given later in this document), but then $m'=|a_2|+|y|<|a_2|+|a_3|=m$, and we assumed that for every $m'<m$ there exist rational solutions of the equation $q'$ in $\mathbb{Q}_p^3$ for all $p$ if and only if there exist rational solutions of $q'$ in $\mathbb{Q}^3$, thus by the bijection above there exist rational solutions of the equation $q$ if and only if there exist nontrivial solutions of $q$ in $\mathbb{Q}_p$ for all $p$, which proves the induction step for $m$.

\section*{$n=4$}
We rearrange $q$ as $q=a_1x_1^2+a_2x_2^2-a_3x_3^2-a_4x_4^2=a_1x_1^2+a_2x_2^2-(a_3x_3^2+a_4x_4^2)$, and we claim that for every prime $p$ there exists $a\in\mathbb{Q}_p$, such that $a_1x_1^2+a_2x_2^2=a_3x_3^2+a_4x_4^2=a$.\\
consider the equation $ay^2-a_1x_1^2+a_2x_2^2=0$, but $a_1x_1^2+a_2x_2^2=a$ for some $v=(v_1,v_2)\in\mathbb{Q}_p^2\implies{a}y^2-a=a(y^2-1)=0$ for $(v_1,v_2,1)\in\mathbb{Q}_p^3$, but this applies for every prime $p$, thus it applies also for $\mathbb{Q}$ itself, as we already observed for $n=3$, that is $ay^2-a_1x_1^2-a_2x_2^2=0$ has a nontrivial solution in $\mathbb{Q}^3$, thus $a_1x_1^2+a_2x_2^2=a$ in $\mathbb{Q}$, and the same exactly applies for $a_3x_3^2+a_4x_4^2=a$ in $\mathbb{Q}$, but then $a_1x_1^2+a_2x_2^2-(a_3x_3^2+a_4x_4^2)=0$ has a nontrivial solution in $\mathbb{Q}$.
\section*{case $n>4$}
We rearrange $q$ as 
\[q=\sum_{i=1}^na_ix_i^2=h-g=a_1x_1^2+a_2x_2^2-(a_3x_3^2+a_4x_4^2+\cdots+a_nx_n^2),\] 
and we assume that the theorem is true for $n-1$ (we have the cases $n=2,3,4$ from above, for this assumption).\\
Suppose that $q(v)=0$ for $v=(v_1,v_2,v_3,v_4,\dots,v_n)\in\mathbb{Q}_p^n$, then there exists some $a\in\mathbb{Q}_p$ such that $h(v')=g(v'')=a_p$, for $v'=(v_1,v_2)$ and $v''=(v_3,v_4,\dots,v_n)$.
We claim (without proof) that the squared $p$-adic rationals form an open set in $\mathbb{Q}_p^{\ast}$, and therefore, due to the approximation theorem, there exists $u=(u_1,u_2)\in\mathbb{Q}^2$ such that if $h(u)=a\in\mathbb{Q}$, then $\frac{a_p}{a}$ is a square in $\mathbb{Q}_p^{\ast}$, but then the quadratic form $g'=ax_{n+1}^2-g$ has a solution\footnote{The quadratic form $g=q-h$ is of dimension $n-2$, and we added $ax_{n+1}^2$.}  $w=(v_3,v_4,\dots,v_n,\sqrt{\frac{a_p}{a}})\in\mathbb{Q}_p^{n-1}$, because $g'(w)=a\frac{a_p}{a}-g(v'')=a_p-a_p=0$, and since $g'$ is of dimension $n-1$, then by the assumption it also has a nontrivial solution $w'=(v'_3,v'_4,\dots,v'_n,v'_{n+1})\in\mathbb{Q}^{n-1}$, but since $g=ax_{n+1}^2-g'$, the vector $w''=(v'_3,v'_4,\dots,v'_n,v'_{n+1},1)$ yields $g(w'')=a-g'(w')=a$, thus the vecotr $u'=(u_1,u_2,v'_1,v'_2,\dots,v'_n,v'_{n+1},1)$ yields $q(u')=h(u)-g(w'')=a-a=0$, which concludes  the proof for all $n>4$, and so we get this result for every $n>1$.
\end{proof}

\label{bijection}
Proposition: Let $a_2,a_3,y,z\in{K}$, such that $z^2-a_2=a_3y$, then
for the following two equations
\begin{equation}
q=x_1^2-a_2x^2-a_3x_3^2=0
\end{equation}
\begin{equation}
q'=x_1^2-a_2x_2^2-yx_3^2=0    
\end{equation}
we have a bijection $f:K^3\rightarrow{K}^3$, such that
\begin{align*}
q(v)&=0\implies{q'}(f(v))=0\\ q'(v)&=0\implies{q}(f^{-1}(v))=0
\end{align*}
\begin{proof}
Let $v=(v_1,v_2,v_3)$ be a solution of  the equation $q$, then we define a map \\$f:K^3\rightarrow{K}^3$
by $(v_1,v_2,v_3)\overset{f}{\mapsto}(a_2v_2+zv_1,zv_2+v_1,a_3v_3)$, and we check that $q'(f(v))=(f(x_1))^2-a_2(f(x_2))^2-y(f(x_3))^2=(a_2v_2+zv_1)^2-a_2(zv_2+v_1)^2-y(a_3v_3)^2=a_2^2v_2^2+2a_2v_2zv_1+z^2v_1^2-a_2z^2v_2^2-2a_2zv_2v_1-a_2v_1^2-ya_3^2v_3^2=(z^2-a_2)v_1^2-(z^2-a_2)a_2v_2^2-(ya_3)a_3v_3^2$, but $z^2-a_2=ya_3$, thus we get that $q'(f(v))=ya_3(v_1^2-a_2v_2^2-a_3v_3^2)=ya_3q(v)=0$.\\
Now we define a map $g:K^3\rightarrow{K}^3$ by $(v_1,v_2,v_3)\overset{g}{\mapsto}(-a_2v_2+zv_1,zv_2-v_1,yv_3)$, and we calculate $(g\circ{f})(v)=g((a_2v_2+zv_1,zv_2+v_1,a_3v_3))=(-a_2(zv_2+v_1)+z(a_2v_2+zv_1),z(zv_2+v_1)-(a_2v_2+zv_1),y(a_3v_3))=(-a_2zv_2-a_2v_1+za_2v_2+z^2v_1,z^2v_2+zv_1-a_2v_2-zv_1,(ya_3)v_3)=((z^2-a_2)v_1,(z^2-a_2)v_2,(ya_3)v_3)=ya_3(v_1,v_2,v_3)$, which shows that $(g\circ{f})(v)=ya_3v$, and we can define $g'=\frac{g}{ya_3}$, thus $(g'\circ{f})(v)=v$, and we take $g$ to be $g'$.\\
\end{proof}
Let $q(v)=0$ in $\mathbb{Q}_p$ for every prime $p$, then there exists $a\in\mathbb{Q}_p$ such that $a_1x_1^2+a_2x_2^2=a$ and $a_3x_3^2+a_4x_4^2=a$.\begin{proof}
The proof is based on a set of corollaries of the basic lemma that every nontrivial isotropic subspace is the orthogonal sum $V=\mathbb{H}^m\perp{U}$, where $U$ is anisotropic and $m=2k$.\\
Consider a quadratic form $q$, where $n=\dim{q}$ is odd, if there exists some $u\in{V}$ such that $q(u)=a$ for some scalar $a\in{K}$, then let $U=Ku<V$ a one-dimensional subspace of $V$, and we construct a subspace $H<V$ by taking a complement to $U$, but then we take the quadratic form  $h=\langle1,-1,1,-1,\dots,1,-1\rangle$ on the $n-1$ basis elements of $H$, thus taking $(1,1,\dots,1)\in{K}^{n-1}$ yields zero, hence the quadratic form $q'=\langle1,-1,\dots,1,-1,a\rangle=h\perp{a}x_{n}^2$ yields $a$ for $(1,1,\dots,1,1)\in{K}^n$, and clearly the quadratic form $q''=q-ax_{n+1}^2$ yields zero for $(1,1,\dots,1,1,1)\in{K}^{n+1}$.\\
Suppose that we start from $q''$ and assume that there exists a $u$ such that $q''(u)=0$, then either $u=(u_1,u_2,\dots,u_n,0)$ such that $q(u_1,u_2,\dots,u_n)=0$, or $u=(u_1,u_2,\dots,u_n,1)$ such that $q(u_1,u_2,\dots,u_n)=a$, thus we get a circular reasoning:\\ $a$ is a value of $q$ $\implies$ $q\sim{h}\perp\langle{a}\rangle$ $\implies$ zero is a value of $q-ax_{n}^2$ $\implies$ $a$ is a value of $q$.\\
Now we assume there exist two quadratic forms $q_1,q_2$ of the same dimension, and we denote the quadratic form of the difference $q=q_1-q_2$. We saw that $a$ is a value of $q_1$ and of $q_1$ if and only if zero is a value of $q_1'=q_1-ax_n^2$ and of $q_2'=q_2-ax_n^2$.

We extend the quadratic form $q$ of the difference to be $q'=q_1\perp-q_2$, and extend the vectors $u$ and $w$ to be $u'=(u_1,u_2,\dots,u_n,0,0,\dots,0)$ and be $w'=(0,0,\dots,0,w_1,u_2,\dots,u_n)$, then if $v=u'+w'$, we have $q'(v)=q_1(u')-q_2(w')=a-a=0$. We observe that if $a=0$, then $q_1(u')=q_2(w')=0$ suggests that every nonzero element in $K$ is a value of $q_1$ and of $q_2$.
\end{proof}
\end{document}